\section{تنظیمات اجرای نهایی ما}

اجرای نهایی روی کارت گرافیکی \lr{NVIDIA Tesla T4} در محیط کگل، پایتون ۳٫۱۲٫۱۳، \lr{PyTorch 2.10.0+cu128} و کودا ۱۲٫۸ انجام شد. برای هر سه مقیاس، همهٔ ۸۰۰ تصویر آموزشی \lr{DIV2K} و جفت‌های رسمی کم‌وضوح دو‌مکعبی به‌کار رفتند. اندازهٔ وصلهٔ پروضوح برای مقیاس‌های ۲، ۳ و ۴ به‌ترتیب ۹۶، ۱۴۴ و ۱۹۲ و اندازهٔ وصلهٔ کم‌وضوح در هر سه حالت ۴۸ بود. اندازهٔ دسته ۳۲، تعداد کارگر بارگذاری داده ۲، بذر تصادفی ۸۱۰۱۰۴۰۳۹ و زیان همان $L_1$ مقاله بود.

برای هر مقیاس دو مرحله اجرا کردیم: مرحلهٔ نخست ۱۰۰ دوره از وزن تصادفی و مرحلهٔ دوم ۱۰۰ دوره با شروع از وزن پایان مرحلهٔ نخست. هر مرحله ۲۵۰۰ به‌روزرسانی و کل آموزش هر مقیاس ۵۰۰۰ به‌روزرسانی داشت. مقاله برای آموزش نخست ۱۰۰۰ دوره را مشخص کرده، اما تعداد دورهٔ بازآموزی مدل اصلی را صریح ننوشته است؛ بنابراین فرض مستند ما این بود که بودجهٔ ۱۰۰‌دوره‌ای انتخاب‌شده برای پروژه، به هر یک از دو مرحله اختصاص یابد. برای جلوگیری از انتخاب نتیجه بر پایهٔ مجموعهٔ آزمون، تمام اعداد اصلی از نقطهٔ وارسی آخرین دوره گرفته شدند، نه بهترین نقطهٔ وارسی بر اساس \lr{Set5}.

جدول \ref{tab:our-training-time} زمان واقعی آموزش را نشان می‌دهد. مجموع زمان پردازش کارت گرافیکی شش اجرا حدود ۹٫۴۱ ساعت بود. اجراهای سه مقیاس مستقل از یکدیگرند و به همین علت امکان اجرای هم‌زمان آن‌ها وجود داشت.

\begin{table}[H]
\centering
\caption{زمان و شمار پارامتر اجرای نهایی؛ زمان‌ها از تاریخچهٔ ذخیره‌شدهٔ خود برنامه استخراج شده‌اند}
\label{tab:our-training-time}
\begin{tabular}{rrrrr}
\toprule
مقیاس & شمار پارامتر & مرحلهٔ نخست (دقیقه) & شروع گرم (دقیقه) & مجموع (دقیقه) \\
\midrule
۲ & ۸۸۱٬۹۱۲ & ۹۵٫۱ & ۹۴٫۴ & ۱۸۹٫۵ \\
۳ & ۸۸۸٬۶۷۷ & ۸۷٫۸ & ۹۶٫۸ & ۱۸۴٫۶ \\
۴ & ۸۹۸٬۱۴۸ & ۹۰٫۷ & ۹۹٫۶ & ۱۹۰٫۳ \\
\bottomrule
\end{tabular}
\end{table}

\section{روند همگرایی}

جدول \ref{tab:our-convergence} زیان نخستین و آخرین دورهٔ هر مرحله را مقایسه می‌کند. در هر سه مقیاس، زیان پایان مرحلهٔ شروع گرم از پایان آموزش نخست کمتر شد. بزرگ‌تر بودن زیان دورهٔ نخست شروع گرم نسبت به آخرین دورهٔ مرحلهٔ نخست تناقض نیست، زیرا ترتیب وصله‌های تصادفی، بهینه‌ساز و حالت‌های مومنتوم در ابتدای مرحلهٔ دوم از نو ساخته شدند.

\begin{table}[H]
\centering
\caption{تغییر میانگین زیان $L_1$ روی دادهٔ کامل \lr{DIV2K}}
\label{tab:our-convergence}
\begin{tabular}{rrrrr}
\toprule
مقیاس & نخستین دورهٔ مرحلهٔ اول & پایان مرحلهٔ اول & نخستین دورهٔ شروع گرم & پایان شروع گرم \\
\midrule
۲ & ۰٫۱۵۷۳۳ & ۰٫۰۱۵۴۲ & ۰٫۰۲۰۹۷ & ۰٫۰۱۴۵۳ \\
۳ & ۰٫۱۹۹۴۳ & ۰٫۰۲۳۳۰ & ۰٫۰۲۹۶۰ & ۰٫۰۲۲۶۵ \\
۴ & ۰٫۱۶۳۵۸ & ۰٫۰۲۸۶۲ & ۰٫۰۳۵۳۶ & ۰٫۰۲۷۰۰ \\
\bottomrule
\end{tabular}
\end{table}

شکل‌های \ref{fig:our-x2-curves} تا \ref{fig:our-x4-curves} هر شش نمودار آموزش را نمایش می‌دهند؛ برای هر مقیاس، نمودار مرحلهٔ نخست و شروع گرم جداگانه آمده است. هر نقطهٔ آبی میانگین زیان تمام ۲۵ دستهٔ آن دوره و نقاط قرمز PSNR پایشی روی \lr{Set5} هستند. افت ناگهانی یا پلهٔ نرخ یادگیری دیده نمی‌شود، زیرا نخستین نصف‌شدن نرخ یادگیری در دورهٔ ۲۰۰ تعریف شده و هر مرحلهٔ ما در دورهٔ ۱۰۰ پایان می‌یابد.

\begin{figure}[H]
    \centering
    \includegraphics[width=0.485\textwidth]{images/training_curve_x2_stage1.png}
    \hfill
    \includegraphics[width=0.485\textwidth]{images/training_curve_x2_warm.png}
    \caption{روند آموزش مدل مقیاس دو روی همهٔ ۸۰۰ تصویر \lr{DIV2K}؛ سمت راست مرحلهٔ ۱۰۰‌دوره‌ای از صفر و سمت چپ مرحلهٔ ۱۰۰‌دوره‌ای شروع گرم است. محور و مقادیر مستقیماً به‌وسیلهٔ برنامه تولید شده‌اند.}
    \label{fig:our-x2-curves}
\end{figure}

\begin{figure}[H]
    \centering
    \includegraphics[width=0.485\textwidth]{images/training_curve_x3_stage1.png}
    \hfill
    \includegraphics[width=0.485\textwidth]{images/training_curve_x3_warm.png}
    \caption{روند کامل آموزش مقیاس سه؛ سمت راست آموزش ۱۰۰‌دوره‌ای از صفر و سمت چپ آموزش ۱۰۰‌دوره‌ای شروع گرم است. هر دو نمودار از تاریخچهٔ ذخیره‌شدهٔ اجرای کگل رسم شده‌اند.}
    \label{fig:our-x3-curves}
\end{figure}

\begin{figure}[H]
    \centering
    \includegraphics[width=0.485\textwidth]{images/training_curve_x4_stage1.png}
    \hfill
    \includegraphics[width=0.485\textwidth]{images/training_curve_x4_warm.png}
    \caption{روند کامل آموزش مقیاس چهار؛ سمت راست آموزش ۱۰۰‌دوره‌ای از صفر و سمت چپ آموزش ۱۰۰‌دوره‌ای شروع گرم است. نوسان PSNR پایشی در این مقیاس علت گزارش‌کردن وزن آخرین دوره به‌جای انتخاب بهترین وزن را روشن می‌کند.}
    \label{fig:our-x4-curves}
\end{figure}

در مقیاس دو، زیان شروع گرم در پایان به ۰٫۰۱۴۵۳ رسید و بیشترین PSNR پایشی نیز در همان انتهای مرحله به‌دست آمد. مقیاس سه روندی مشابه داشت و زیان نهایی آن ۰٫۰۲۲۶۵ شد. در مقیاس چهار، زیان تا ۰٫۰۲۷۰۰ کاهش یافت، اما PSNR پایشی به‌علت حساسیت بیشتر بازسازی چهاربرابری نوسان داشت؛ بهترین مقدار پایشی ۳۰٫۹۶۷ و مقدار آخرین دوره ۳۰٫۷۷۹ دسی‌بل بود. ما مقدار دوم را در تمام جدول‌های اصلی به‌کار بردیم تا قاعدهٔ انتخاب وزن برای هر سه مقیاس یکسان بماند.

شکل \ref{fig:combined-training-loss} همین روند را به‌صورت فشرده نمایش می‌دهد. در هر سه مقیاس، منحنی مرحلهٔ شروع گرم از مقدار زیان پایین‌تری آغاز می‌شود و در پایان نیز از مرحلهٔ نخست کمتر است.

\begin{figure}[H]
    \centering
    \includegraphics[width=\textwidth]{images/training_loss_curves.png}
    \caption{مقایسهٔ زیان آموزش در دو مرحلهٔ ۱۰۰‌دوره‌ای برای هر سه مقیاس. منحنی‌ها مستقیماً از تاریخچه‌های ذخیره‌شدهٔ اجرای نهایی رسم شده‌اند.}
    \label{fig:combined-training-loss}
\end{figure}

\section{نتایج نهایی روی پنج مجموعهٔ معیار}

برای اینکه اثر هر مرحله مستقل دیده شود، جدول \ref{tab:our-stage-one-results} همهٔ نتایج پایان آموزش نخست را گزارش می‌کند. سپس جدول \ref{tab:our-full-results} نتیجهٔ پایان دورهٔ ۱۰۰ شروع گرم را برای همان ۱۵ ترکیب مقیاس و مجموعه می‌آورد. سنجه‌ها روی کانال روشنایی و پس از حذف حاشیه‌ای به عرض ضریب مقیاس محاسبه شده‌اند. هر مقدار میانگین همهٔ تصاویر آن مجموعه است، نه نتیجهٔ یک نمونهٔ منتخب: ۵ تصویر \lr{Set5}، ۱۴ تصویر \lr{Set14}، ۱۰۰ تصویر \lr{B100}، ۱۰۰ تصویر \lr{Urban100} و ۱۰۹ تصویر \lr{Manga109}؛ در مجموع ۳۲۸ تصویر در هر اجرا. نتیجهٔ تصویر‌به‌تصویر شش اجرا نیز در ۳۰ فایل CSV داخل پوشه‌های \lr{outputs} تحویل داده شده است. این ۱۹۶۸ ردیف خام مبنای میانگین‌های گزارش‌اند، اما برای خوانایی به‌صورت کامل داخل متن چاپ نشده‌اند.

\begin{table}[H]
\centering
\caption{همهٔ نتایج پایان مرحلهٔ نخست ۱۰۰‌دوره‌ای پیش از شروع گرم؛ هر خانه به‌ترتیب PSNR برحسب دسی‌بل و SSIM را نشان می‌دهد}
\label{tab:our-stage-one-results}
\resizebox{\textwidth}{!}{%
\begin{tabular}{rrrrrr}
\toprule
مقیاس & \lr{Set5} & \lr{Set14} & \lr{B100} & \lr{Urban100} & \lr{Manga109} \\
\midrule
۲ & ۳۶٫۷۵۹ / ۰٫۹۵۴۳ & ۳۲٫۴۸۲ / ۰٫۹۰۶۶ & ۳۱٫۳۸۷ / ۰٫۸۸۸۰ & ۲۹٫۶۶۸ / ۰٫۸۹۷۳ & ۳۶٫۰۷۵ / ۰٫۹۶۵۳ \\
۳ & ۳۲٫۴۶۷ / ۰٫۹۰۲۳ & ۲۹٫۱۱۴ / ۰٫۸۱۶۵ & ۲۸٫۲۶۰ / ۰٫۷۸۲۹ & ۲۵٫۹۳۹ / ۰٫۷۹۱۷ & ۲۹٫۸۶۱ / ۰٫۸۹۹۰ \\
۴ & ۳۰٫۲۵۴ / ۰٫۸۵۷۳ & ۲۷٫۳۵۰ / ۰٫۷۵۰۸ & ۲۶٫۷۹۳ / ۰٫۷۱۱۰ & ۲۴٫۳۲۳ / ۰٫۷۱۸۴ & ۲۷٫۰۷۶ / ۰٫۸۴۳۵ \\
\bottomrule
\end{tabular}}
\end{table}

مرحلهٔ نخست در همهٔ مجموعه‌ها از دو‌مکعبی بهتر شد، ولی فاصلهٔ آن با خروجی شروع گرم روی مجموعه‌های بافت‌دار بیشتر است. برای نمونه، شروع گرم PSNR مقیاس سه را روی \lr{Manga109} از ۲۹٫۸۶۱ به ۳۱٫۲۶۹ رساند؛ یعنی افزایش ۱٫۴۰۸ دسی‌بل بدون تغییر معماری یا شمار پارامتر.

\begin{table}[H]
\centering
\caption{همهٔ نتایج پایان مرحلهٔ شروع گرم در مقایسه با درون‌یابی دو‌مکعبی و STSN گزارش‌شده در مقاله؛ هر جفت ستون شامل PSNR برحسب دسی‌بل و SSIM است}
\label{tab:our-full-results}
\resizebox{\textwidth}{!}{%
\begin{tabular}{clrrrrrr}
\toprule
مقیاس & مجموعه & \multicolumn{2}{c}{مدل ما} & \multicolumn{2}{c}{دو‌مکعبی} & \multicolumn{2}{c}{مقاله} \\
 & & PSNR & SSIM & PSNR & SSIM & PSNR & SSIM \\
\midrule
۲ & \lr{Set5}     & ۳۷٫۱۶۳ & ۰٫۹۵۷۵ & ۳۳٫۹۶۰ & ۰٫۹۳۳۶ & ۳۸٫۱۹ & ۰٫۹۶۱۱ \\
۲ & \lr{Set14}    & ۳۲٫۸۲۶ & ۰٫۹۱۱۵ & ۳۰٫۵۴۴ & ۰٫۸۷۵۳ & ۳۳٫۷۸ & ۰٫۹۱۹۹ \\
۲ & \lr{B100}     & ۳۱٫۶۸۵ & ۰٫۸۹۳۶ & ۲۹٫۷۴۳ & ۰٫۸۴۹۷ & ۳۲٫۳۰ & ۰٫۹۰۱۳ \\
۲ & \lr{Urban100} & ۳۰٫۳۴۳ & ۰٫۹۰۹۳ & ۲۷٫۰۷۳ & ۰٫۸۴۵۸ & ۳۲٫۶۸ & ۰٫۹۳۳۶ \\
۲ & \lr{Manga109} & ۳۶٫۸۶۷ & ۰٫۹۷۲۰ & ۳۱٫۲۵۲ & ۰٫۹۳۸۶ & ۳۹٫۱۳ & ۰٫۹۷۷۸ \\
\midrule
۳ & \lr{Set5}     & ۳۳٫۱۷۹ & ۰٫۹۱۶۶ & ۳۰٫۶۱۳ & ۰٫۸۷۲۱ & ۳۴٫۶۲ & ۰٫۹۲۹۲ \\
۳ & \lr{Set14}    & ۲۹٫۴۶۶ & ۰٫۸۲۷۷ & ۲۷٫۷۸۷ & ۰٫۷۸۱۲ & ۳۰٫۵۴ & ۰٫۸۴۶۶ \\
۳ & \lr{B100}     & ۲۸٫۵۷۰ & ۰٫۷۹۳۲ & ۲۷٫۳۱۹ & ۰٫۷۴۴۳ & ۲۹٫۲۲ & ۰٫۸۰۹۰ \\
۳ & \lr{Urban100} & ۲۶٫۶۵۵ & ۰٫۸۱۶۴ & ۲۴٫۵۸۲ & ۰٫۷۳۹۸ & ۲۸٫۵۹ & ۰٫۸۶۲۱ \\
۳ & \lr{Manga109} & ۳۱٫۲۶۹ & ۰٫۹۲۳۶ & ۲۷٫۱۹۸ & ۰٫۸۶۰۳ & ۳۴٫۱۱ & ۰٫۹۴۸۰ \\
\midrule
۴ & \lr{Set5}     & ۳۰٫۷۷۹ & ۰٫۸۷۲۲ & ۲۸٫۶۰۱ & ۰٫۸۱۴۲ & ۳۲٫۴۶ & ۰٫۸۹۸۲ \\
۴ & \lr{Set14}    & ۲۷٫۷۲۳ & ۰٫۷۶۲۴ & ۲۶٫۲۱۰ & ۰٫۷۰۹۰ & ۲۸٫۷۶ & ۰٫۷۸۶۰ \\
۴ & \lr{B100}     & ۲۷٫۰۵۸ & ۰٫۷۱۹۷ & ۲۶٫۰۵۱ & ۰٫۶۷۲۲ & ۲۷٫۶۸ & ۰٫۷۴۰۵ \\
۴ & \lr{Urban100} & ۲۴٫۷۸۱ & ۰٫۷۳۹۳ & ۲۳٫۲۴۰ & ۰٫۶۶۱۶ & ۲۶٫۳۹ & ۰٫۷۹۷۱ \\
۴ & \lr{Manga109} & ۲۸٫۱۶۸ & ۰٫۸۶۹۷ & ۲۵٫۰۷۲ & ۰٫۷۹۰۶ & ۳۰٫۹۳ & ۰٫۹۱۴۲ \\
\bottomrule
\end{tabular}}
\end{table}

شکل \ref{fig:benchmark-psnr} مقایسهٔ جدول را به‌صورت بصری خلاصه می‌کند. در همهٔ مجموعه‌ها و هر سه مقیاس، خروجی نهایی STSN از درون‌یابی دو‌مکعبی PSNR بالاتری دارد.

\begin{figure}[H]
    \centering
    \includegraphics[width=\textwidth]{images/benchmark_psnr_vs_bicubic.png}
    \caption{مقایسهٔ PSNR خروجی نهایی STSN پس از شروع گرم با خط پایهٔ دو‌مکعبی روی پنج مجموعهٔ معیار. هر پنل یک ضریب بزرگ‌نمایی را نشان می‌دهد.}
    \label{fig:benchmark-psnr}
\end{figure}

در همهٔ ۱۵ حالت، مدل آموزش‌دیده از درون‌یابی دو‌مکعبی بهتر است. بیشترین بهبود در \lr{Manga109} رخ داد: ۵٫۶۱۵، ۴٫۰۷۱ و ۳٫۰۹۶ دسی‌بل برای مقیاس‌های ۲، ۳ و ۴. این رفتار با ماهیت این مجموعه سازگار است، زیرا خطوط و بافت‌های تکراری تصویرهای کمیک از یادگیری ساختاری بیش از هموارسازی ثابت سود می‌برند. کمترین بهبود مربوط به \lr{B100} بود، ولی حتی آن‌جا نیز برای هر سه مقیاس بیش از یک دسی‌بل افزایش به‌دست آمد.

\subsection{تحلیل مجموعه‌به‌مجموعه}

در \lr{Set5} که کوچک‌ترین و معمولاً ساده‌ترین معیار است، خروجی نهایی به‌ترتیب ۳۷٫۱۶۳، ۳۳٫۱۷۹ و ۳۰٫۷۷۹ دسی‌بل شد. SSIM نیز از ۰٫۹۵۷۵ در مقیاس دو به ۰٫۸۷۲۲ در مقیاس چهار کاهش یافت. این افت با افزایش ضریب بزرگ‌نمایی طبیعی است، زیرا در مقیاس چهار تعداد بیشتری از جزئیات باید از ورودی کوچک‌تر تخمین زده شوند.

در \lr{Set14} تنوع محتوایی بیشتر است و نتیجهٔ سه مقیاس به ۳۲٫۸۲۶، ۲۹٫۴۶۶ و ۲۷٫۷۲۳ دسی‌بل رسید. بهبود نسبت به دو‌مکعبی بین ۱٫۵۱۳ تا ۲٫۲۸۳ دسی‌بل بود؛ بنابراین شبکه فقط روی مجموعهٔ کوچک \lr{Set5} بهتر نشده و روی تصاویر متنوع‌تر نیز مزیت پایدار دارد.

در \lr{B100} کمترین افزایش PSNR نسبت به دو‌مکعبی، یعنی ۱٫۹۴۲، ۱٫۲۵۱ و ۱٫۰۰۷ دسی‌بل، دیده شد. این مجموعه بیشتر شامل تصاویر طبیعی است و بافت‌های نامنظم آن نسبت به خطوط تکراری شهری و کمیک، سود کمتری از بازسازی ساختاری شبکه می‌برند. در عوض فاصلهٔ ما با مقاله در همین مجموعه از همه کمتر بود: ۰٫۶۱۵، ۰٫۶۵۰ و ۰٫۶۲۲ دسی‌بل.

در \lr{Urban100} شروع گرم نقش مهم‌تری داشت و برای مقیاس‌های ۲، ۳ و ۴ به‌ترتیب ۰٫۶۷۵، ۰٫۷۱۶ و ۰٫۴۵۸ دسی‌بل افزود. خروجی نهایی ۳۰٫۳۴۳، ۲۶٫۶۵۵ و ۲۴٫۷۸۱ دسی‌بل بود. خطوط نما، پنجره‌ها و الگوهای تکراری ساختمان‌ها با میدان دید ۱۱ در ۱۱ پیمانه‌سازی و توجه فضایی معماری سازگارند، ولی فاصلهٔ باقی‌مانده با مقاله نشان می‌دهد یادگیری همین جزئیات به آموزش طولانی‌تر حساس است.

در \lr{Manga109} هم PSNR و هم SSIM بیشترین سود را از مدل گرفتند. خروجی نهایی به ۳۶٫۸۶۷/۰٫۹۷۲۰، ۳۱٫۲۶۹/۰٫۹۲۳۶ و ۲۸٫۱۶۸/۰٫۸۶۹۷ رسید. افزایش شروع گرم در مقیاس‌های سه و چهار به‌ترتیب ۱٫۴۰۸ و ۱٫۰۹۲ دسی‌بل بود؛ این دو مقدار بزرگ‌ترین اثر شروع گرم در کل جدول‌اند. در نتیجه، آموزش دوم بیش از همه برای بازسازی خطوط باریک و بافت‌های منظم کمیک مؤثر بوده است.

جدول \ref{tab:our-psnr-differences} دو مقایسهٔ مهم را جدا می‌کند. ستون نخست نشان می‌دهد شروع گرم در همهٔ حالت‌ها مفید بوده است. ستون دوم فاصلهٔ خروجی نهایی با خط پایه را نشان می‌دهد. ستون سوم اختلاف با مقاله است؛ منفی بودن آن یعنی اجرای ۲۰۰‌دوره‌ای ما هنوز به آموزش طولانی مقاله نرسیده است.

\begin{table}[H]
\centering
\caption{اختلاف PSNR خروجی نهایی برحسب دسی‌بل؛ «مرحلهٔ اول» خروجی ۱۰۰‌دوره‌ای پیش از شروع گرم است}
\label{tab:our-psnr-differences}
\resizebox{0.88\textwidth}{!}{%
\begin{tabular}{clrrr}
\toprule
مقیاس & مجموعه & نسبت به مرحلهٔ اول & نسبت به دو‌مکعبی & نسبت به مقاله \\
\midrule
۲ & \lr{Set5}     & +۰٫۴۰۴ & +۳٫۲۰۳ & -۱٫۰۲۷ \\
۲ & \lr{Set14}    & +۰٫۳۴۵ & +۲٫۲۸۳ & -۰٫۹۵۴ \\
۲ & \lr{B100}     & +۰٫۲۹۸ & +۱٫۹۴۲ & -۰٫۶۱۵ \\
۲ & \lr{Urban100} & +۰٫۶۷۵ & +۳٫۲۷۱ & -۲٫۳۳۷ \\
۲ & \lr{Manga109} & +۰٫۷۹۲ & +۵٫۶۱۵ & -۲٫۲۶۳ \\
۳ & \lr{Set5}     & +۰٫۷۱۳ & +۲٫۵۶۶ & -۱٫۴۴۱ \\
۳ & \lr{Set14}    & +۰٫۳۵۱ & +۱٫۶۷۹ & -۱٫۰۷۴ \\
۳ & \lr{B100}     & +۰٫۳۱۰ & +۱٫۲۵۱ & -۰٫۶۵۰ \\
۳ & \lr{Urban100} & +۰٫۷۱۶ & +۲٫۰۷۳ & -۱٫۹۳۵ \\
۳ & \lr{Manga109} & +۱٫۴۰۸ & +۴٫۰۷۱ & -۲٫۸۴۱ \\
۴ & \lr{Set5}     & +۰٫۵۲۵ & +۲٫۱۷۷ & -۱٫۶۸۱ \\
۴ & \lr{Set14}    & +۰٫۳۷۳ & +۱٫۵۱۳ & -۱٫۰۳۷ \\
۴ & \lr{B100}     & +۰٫۲۶۵ & +۱٫۰۰۷ & -۰٫۶۲۲ \\
۴ & \lr{Urban100} & +۰٫۴۵۸ & +۱٫۵۴۰ & -۱٫۶۰۹ \\
۴ & \lr{Manga109} & +۱٫۰۹۲ & +۳٫۰۹۶ & -۲٫۷۶۲ \\
\bottomrule
\end{tabular}}
\end{table}

شکل \ref{fig:warm-start-gain} افزایش PSNR مرحلهٔ شروع گرم نسبت به مرحلهٔ نخست را برای همهٔ ۱۵ حالت نشان می‌دهد. همهٔ میله‌ها مثبت‌اند؛ بیشترین افزایش مربوط به \lr{Manga109} در مقیاس سه است.

\begin{figure}[H]
    \centering
    \includegraphics[width=0.9\textwidth]{images/warm_start_psnr_gain.png}
    \caption{افزایش PSNR بر اثر مرحلهٔ ۱۰۰‌دوره‌ای شروع گرم نسبت به پایان مرحلهٔ نخست، برای هر مجموعه و مقیاس.}
    \label{fig:warm-start-gain}
\end{figure}

فاصله با مقاله روی \lr{B100} کمتر و روی مجموعه‌های ساختاری \lr{Urban100} و \lr{Manga109} بیشتر است. برداشت ما این است که بافت‌ها و وابستگی‌های پیچیده‌تر برای استفادهٔ کامل از میدان دید بزرگ شبکه به تعداد به‌روزرسانی بیشتری نیاز دارند. این یک استنباط از الگوی جدول است، نه آزمایش حذف مستقل. همچنین در مقیاس چهار، بهترین نقطهٔ وارسی پایشی \lr{Set5} به ۳۰٫۹۶۷ دسی‌بل رسید، اما برای رعایت قاعدهٔ یکسان و پرهیز از انتخاب روی آزمون، مقدار پایان دوره یعنی ۳۰٫۷۷۹ گزارش شد.

\section{زمان استنتاج و بررسی کیفی}

جدول \ref{tab:our-inference-time} میانگین زمان اجرای یک تصویر را روی همان کارت T4 نشان می‌دهد. زمان به اندازه و مقیاس تصویر وابسته است؛ ازاین‌رو مقایسهٔ مستقیم اعداد مجموعه‌های متفاوت یا مقایسه با سخت‌افزار مقاله منصفانه نیست. الگوی کلی درست است: با بزرگ‌تر شدن ورودی کم‌وضوح، زمان افزایش می‌یابد و \lr{Manga109} و \lr{Urban100} سنگین‌ترین مجموعه‌ها هستند.

\begin{table}[H]
\centering
\caption{میانگین زمان استنتاج هر تصویر روی کارت \lr{Tesla T4}، برحسب میلی‌ثانیه}
\label{tab:our-inference-time}
\begin{tabular}{rrrrrr}
\toprule
مقیاس & \lr{Set5} & \lr{Set14} & \lr{B100} & \lr{Urban100} & \lr{Manga109} \\
\midrule
۲ & ۶۰٫۵ & ۱۱۴٫۶ & ۷۳٫۰ & ۴۰۱٫۴ & ۵۱۶٫۷ \\
۳ & ۳۰٫۸ & ۶۰٫۶ & ۳۹٫۸ & ۱۸۲٫۲ & ۲۲۷٫۵ \\
۴ & ۲۰٫۶ & ۴۷٫۰ & ۲۸٫۵ & ۹۸٫۰ & ۱۲۸٫۴ \\
\bottomrule
\end{tabular}
\end{table}

\subsection{نمونه‌های کیفی در حالت‌های مختلف}

برای اینکه بررسی بصری به یک تصویر شهری محدود نماند، هفت نمونه از پنج مجموعه و هر سه ضریب بزرگ‌نمایی در شکل‌های \ref{fig:qual-set5-baby-x2} تا \ref{fig:qual-urban-x4} آورده‌ایم. هر چهار پنل مستقیماً از آرایه‌های ورودی، درون‌یابی، خروجی شبکه و مرجع ساخته شده‌اند و تصویر صفحه نیستند. این نمونه‌ها برای پوشش‌دادن چهره و بافت پارچه، بافت طبیعی، خطوط معماری، تصویر کمیک و بزرگ‌نمایی دشوار چهاربرابری انتخاب شده‌اند. همهٔ تصاویر مقایسه‌ای تولیدشده، از جمله نمونه‌های دیگری که در متن جا نشده‌اند، در پوشهٔ \lr{benchmarks} هر اجرا موجودند.

\begin{table}[H]
\centering
\caption{سنجه‌های عددی نمونه‌های کیفی نمایش‌داده‌شده؛ اعداد از ردیف همان تصویر در فایل‌های CSV گرفته شده‌اند}
\label{tab:qualitative-samples}
\resizebox{0.95\textwidth}{!}{%
\begin{tabular}{cllrrrr}
\toprule
مقیاس & مجموعه & تصویر & PSNR مدل & SSIM مدل & PSNR دو‌مکعبی & SSIM دو‌مکعبی \\
\midrule
۲ & \lr{Set5} & \lr{baby} & ۳۸٫۶۷۵ & ۰٫۹۶۶۴ & ۳۷٫۲۳۶ & ۰٫۹۵۴۹ \\
۲ & \lr{Set14} & \lr{barbara} & ۲۸٫۳۴۲ & ۰٫۸۷۶۵ & ۲۸٫۱۰۷ & ۰٫۸۴۶۷ \\
۲ & \lr{Urban100} & \lr{img\_001} & ۳۱٫۷۲۹ & ۰٫۸۹۴۹ & ۲۸٫۳۰۶ & ۰٫۸۲۶۱ \\
۳ & \lr{B100} & \lr{101085} & ۲۴٫۷۹۳ & ۰٫۶۵۲۳ & ۲۴٫۰۴۱ & ۰٫۵۹۰۷ \\
۳ & \lr{Manga109} & \lr{ARMS} & ۲۹٫۵۱۸ & ۰٫۸۵۲۸ & ۲۷٫۰۷۶ & ۰٫۷۷۵۴ \\
۴ & \lr{Set5} & \lr{bird} & ۳۳٫۰۷۹ & ۰٫۹۲۱۶ & ۳۰٫۴۴۰ & ۰٫۸۷۷۶ \\
۴ & \lr{Urban100} & \lr{img\_002} & ۲۵٫۶۱۰ & ۰٫۷۷۰۹ & ۲۴٫۲۷۷ & ۰٫۶۸۲۸ \\
\bottomrule
\end{tabular}}
\end{table}

در نمونهٔ \lr{baby}، مدل بافت کلاه، مرز چشم‌ها و خطوط ظریف صورت را بدون ایجاد لبهٔ شدید حفظ کرده و ۱٫۴۳۹ دسی‌بل از دو‌مکعبی بهتر است. این تصویر نمونه‌ای از ناحیه‌های نرم همراه با چند بافت ریز است.

\begin{figure}[H]
    \centering
    \includegraphics[width=\textwidth]{images/qualitative_set5_baby_x2.png}
    \caption{نمونهٔ \lr{baby} از \lr{Set5} در مقیاس دو پس از شروع گرم؛ از راست به چپ: مرجع پروضوح، STSN، دو‌مکعبی و ورودی کم‌وضوح.}
    \label{fig:qual-set5-baby-x2}
\end{figure}

نمونهٔ \lr{barbara} به‌خاطر پارچه‌های راه‌راه و بافت‌های بسامدبالا دشوار است. افزایش PSNR فقط ۰٫۲۳۴ دسی‌بل است، اما SSIM از ۰٫۸۴۶۷ به ۰٫۸۷۶۵ رسیده؛ یعنی بهبود ساختار محلی از کاهش خطای پیکسلی مشهودتر بوده است.

\begin{figure}[H]
    \centering
    \includegraphics[width=\textwidth]{images/qualitative_set14_barbara_x2.png}
    \caption{نمونهٔ \lr{barbara} از \lr{Set14} در مقیاس دو؛ بافت پارچه و خطوط ریز میز، تفاوت بازسازی ساختاری را آشکار می‌کنند. ترتیب پنل‌ها از راست: مرجع، STSN، دو‌مکعبی و ورودی است.}
    \label{fig:qual-set14-barbara-x2}
\end{figure}

در شکل \ref{fig:our-qualitative}، تکرار پنجره‌ها و لبه‌های ساختمان نسبت به دو‌مکعبی منظم‌تر شده و افزایش ۳٫۴۲۳ دسی‌بلی با این تفاوت بصری سازگار است. بااین‌حال، در جزئیات بسیار ریز هنوز اختلاف با مرجع دیده می‌شود.

\begin{figure}[H]
    \centering
    \includegraphics[width=\textwidth]{images/comparison_x2_urban100_warm.png}
    \caption{مقایسهٔ کیفی یک تصویر \lr{Urban100} در مقیاس دو پس از شروع گرم؛ از راست به چپ: مرجع پروضوح، خروجی STSN، درون‌یابی دو‌مکعبی و ورودی کم‌وضوح. شکل خروجی مستقیم برنامه است.}
    \label{fig:our-qualitative}
\end{figure}

نمونهٔ \lr{101085} از \lr{B100} بافت طبیعی و نامنظم دارد. در این‌جا افزایش PSNR برابر ۰٫۷۵۲ دسی‌بل است، ولی SSIM حدود ۰٫۰۶۲ افزایش یافته است. مدل مرز پیکره‌ها را واضح‌تر می‌کند، اما نبود الگوی تکراری منظم باعث می‌شود سود آن از نمونه‌های شهری و کمیک کمتر باشد.

\begin{figure}[H]
    \centering
    \includegraphics[width=\textwidth]{images/qualitative_b100_101085_x3.png}
    \caption{نمونهٔ \lr{101085} از \lr{B100} در مقیاس سه؛ از راست به چپ: مرجع، STSN، دو‌مکعبی و ورودی کم‌وضوح. این نمونه رفتار مدل روی بافت طبیعی نامنظم را نشان می‌دهد.}
    \label{fig:qual-b100-x3}
\end{figure}

در تصویر \lr{ARMS} خطوط باریک نقشه، نوشته‌ها و مرزهای ترسیمی فراوان‌اند. STSN نسبت به دو‌مکعبی ۲٫۴۴۲ دسی‌بل و حدود ۰٫۰۷۷ SSIM بهتر است. این نتیجه با بهبود بالای میانگین \lr{Manga109} هم‌جهت است و نشان می‌دهد مدل در بازسازی ساختارهای خطی و تکرارشونده موفق‌تر است.

\begin{figure}[H]
    \centering
    \includegraphics[width=\textwidth]{images/qualitative_manga109_arms_x3.png}
    \caption{نمونهٔ \lr{ARMS} از \lr{Manga109} در مقیاس سه؛ خروجی STSN خطوط ترسیمی و نوشته‌های ریز را نسبت به دو‌مکعبی واضح‌تر نگه می‌دارد.}
    \label{fig:qual-manga-x3}
\end{figure}

برای نمایش دشواری مقیاس چهار، دو محتوای متفاوت انتخاب شد. در تصویر \lr{bird}، مرز منقار و چشم و شاخه‌های باریک واضح‌تر شده‌اند و افزایش PSNR به ۲٫۶۳۹ دسی‌بل می‌رسد؛ بااین‌حال، فاصله با مرجع هنوز بیشتر از نمونه‌های مقیاس دو است.

\begin{figure}[H]
    \centering
    \includegraphics[width=\textwidth]{images/qualitative_set5_bird_x4.png}
    \caption{نمونهٔ \lr{bird} از \lr{Set5} در بزرگ‌نمایی چهاربرابری؛ این حالت اثر کاهش شدید اطلاعات ورودی بر مرزهای طبیعی را نشان می‌دهد. ترتیب پنل‌ها از راست: مرجع، STSN، دو‌مکعبی و ورودی است.}
    \label{fig:qual-set5-bird-x4}
\end{figure}

نمونهٔ شهری مقیاس چهار شامل خط‌های مورب و منحنی‌های تکرارشوندهٔ سازه است. مدل ۱٫۳۳۳ دسی‌بل و حدود ۰٫۰۸۸ SSIM بهتر از دو‌مکعبی عمل می‌کند و پیوستگی تیرها را بهتر نگه می‌دارد، ولی در بخش‌های بسیار متراکم بخشی از خطوط ظریف همچنان نرم می‌شوند.

\begin{figure}[H]
    \centering
    \includegraphics[width=\textwidth]{images/qualitative_urban100_img002_x4.png}
    \caption{نمونهٔ \lr{img\_002} از \lr{Urban100} در مقیاس چهار؛ خطوط مورب و قوس‌های تکراری سازه برای مقایسهٔ بازسازی هندسی مناسب‌اند.}
    \label{fig:qual-urban-x4}
\end{figure}
